\originalchapter{谱线、相关函数与精密光谱}
\originalsection{从有限相干时间到谱线}
弱外场的激发幅度为$A(t)\propto\int_0^t x(t')\e^{\ii\Delta t'}\dd t'$. 概率平均是一个双积分. 若过程平稳, 定义$C(\tau)=\langle x^*(t)x(t+\tau)\rangle$, 把积分变量换成时差与共同时间, 在$t$远长于相关时间时得
\[\frac{\langle|A(t)|^2\rangle}{t}\ \propto\ \int_{-\infty}^{\infty}C(\tau)\e^{\ii\Delta\tau}\dd\tau.\]
短时相位尚未随机化, 概率二次增长; 长时各相干区间随机相加, 概率线性增长. 速率方程的适用范围因此可由相关时间解释.

若$C(\tau)=C(0)\e^{-|\tau|/T_2}$, 积分为$2C(0)T_2/[1+(\Delta T_2)^2]$, 对应Lorentz线, 全宽$2/T_2$. 静态Gaussian频率分布方差$\sigma_\omega^2$给$C(\tau)=C(0)\e^{-\sigma_\omega^2\tau^2/2}$与Gaussian谱. 独立展宽机制的相关函数相乘, 频谱就卷积; Lorentz宽度相加, Gaussian方差相加, 两者混合为Voigt线而非任选一个宽度相加.

\originalsection{Doppler、碰撞与运动变窄}
一维热速度方差为$k_BT/m$, 运动原子看到失谐$\Delta-kv_z$. 因此Doppler角频率标准差$k\sqrt{k_BT/m}$, Gaussian全宽为$2\sqrt{2\ln2}\,k\sqrt{k_BT/m}$. 两束反向的光做饱和谱时, 只有特定速度组同时与两束共振; 在零速度组产生窄Lamb凹陷. 另有不同跃迁共享速度组的交叉峰, 不能把所有窄峰都认作真实能级差.

若频率噪声$\delta\omega$是Gaussian且$\langle\delta\omega(0)\delta\omega(t)\rangle=\sigma^2\e^{-|t|/\tau_c}$, 累积相位$\phi(t)=\int_0^t\delta\omega(s)\dd s$有
\[\langle\e^{\ii\phi(t)}\rangle=\exp\{-\sigma^2\tau_c[t-\tau_c(1-\e^{-t/\tau_c})]\}.\]
这由$\langle\phi^2\rangle=2\int_0^t(t-s)\sigma^2\e^{-s/\tau_c}\dd s$及Gaussian特征函数得到. 慢噪声给Gaussian初始衰减; 快噪声给衰减率$\sigma^2\tau_c$, 越快反而越窄. 原子扩散满足$\langle\e^{\ii\mathbf k\cdot[\mathbf r(t)-\mathbf r(0)]}\rangle=\e^{-Dk^2t}$, 所以Dicke变窄的角频率全宽为$2Dk^2$. 有些课件把半宽记为“linewidth”; 比较数值必须说明半宽还是全宽.

\originalsection{双光子与相位匹配}
二阶跃迁振幅有两个时间顺序. 对吸收$\omega_1,\omega_2$的过程, 有效耦合含
\[\sum_n\left[\frac{(\mathbf d_{fn}\cdot\mathbf E_2)(\mathbf d_{ni}\cdot\mathbf E_1)}{E_i+\hbar\omega_1-E_n}+\frac{(\mathbf d_{fn}\cdot\mathbf E_1)(\mathbf d_{ni}\cdot\mathbf E_2)}{E_i+\hbar\omega_2-E_n}\right].\]
非共振中间态是虚态, 不要求真的存在中间态布居. 两束反向且频率接近时, 净波矢$\mathbf k_1+\mathbf k_2$很小, 一阶Doppler位移受抑制. 同向光则保留两倍单光子Doppler灵敏度. $S\to S$电偶极单光子禁戒, 但两个奇宇称偶极作用的乘积为偶宇称, 因而两光子允许.

两束频差作用于自由原子形成运动光栅. 若转移动量$\hbar\mathbf q$, 能量守恒要求$\hbar\delta=[(\mathbf p+\hbar\mathbf q)^2-p^2]/(2m)$, 即$\delta=\mathbf q\cdot\mathbf v+\hbar q^2/(2m)$. 这就是Bragg共振, 包括Doppler项与反冲项, 不能只用固定原子二能级模型而忽略运动.
